% from ega4-21a-fr.tex
\oldpage[IV-4]{260}

On en conclut que $\sh{J}_X(D)$, pour tout diviseur $D\in\operatorname{Div}(X)$, est l'Idéal fractionnaire inversible défini de la façon suivante : pour tout ouvert $U$ de $X$ tel que $D|U=\operatorname{div}_U(f)$, où $f\in\Gamma(U,\sh{M}_X^*)$, $\sh{J}_X(D)|U$ est l'Idéal fractionnaire inversible $\sh{O}_U.f$.
On a donc, en vertu de \sref{IV.21.1.6}, pour toute fonction méromorphe régulière $f\in\Gamma(X,\sh{M}_X^*)$, la relation
\[
\label{IV.21.2.5.4}
  f\in\Gamma(X,\sh{J}_X(D)) \iff \operatorname{div}(f)\geq D.
  \tag{21.2.5.4}
\]
\end{env}

\begin{proposition}[21.2.6]
\label{IV.21.2.6}
L'homomorphisme $I:\sh{D}iv_X\to\sh{J}d.inv_X$ est bijectif.
\end{proposition}

\begin{proof}
On définit en effet un homomorphisme $I_X'$ de $\operatorname{Id.inv}(X)$ dans $\operatorname{Div}(X)$ en faisant correspondre à tout Idéal fractionnaire inversible $\sh{J}$ sur $X$ le diviseur $I_X'(\sh{J})$ défini de la façon suivante : pour tout ouvert $U$ de $X$ tel que $\sh{J}|U=\sh{O}_U.f$, où $f\in\Gamma(U,\sh{M}_X^*)$ \sref{IV.21.2.2}, on prend $I_X'(\sh{J})|U=\operatorname{div}_U(f)$;
en vertu de la remarque suivant \sref{IV.21.2.2}, cette définition est indépendante du générateur $f$ choisi dans $\sh{J}|U$, et détermine bien un diviseur sur $X$.
En outre, cette définition montre aussitôt que les homomorphismes $\sh{J}_X$ et $I_X'$ sont réciproques l'un de l'autre.
Remplaçant $X$ par un ouvert quelconque $U$, on en déduit la définition de l'isomorphisme $I':\sh{J}d.inv_X\to\sh{D}iv_X$ réciproque de $I$, d'où la proposition.
On posera $I_X'(\sh{J})=\operatorname{div}(\sh{J})$, de sorte que l'on a, pour toute fonction méromorphe régulière $f$ sur $X$
\[
\label{IV.21.2.6.1}
  \operatorname{div}(\sh{O}_X.f)=\operatorname{div}(f).
  \tag{21.2.6.1}
\]
\end{proof}

\begin{env}[21.2.7]
\EGAReaderTarget{ega4.c04.statement.21-2-7}{IV.21.2.7}
On identifiera souvent les faisceaux $\sh{D}iv_X$ et $\sh{J}d.inv_X$ (resp. les groupes $\operatorname{Div}(X)$ et $\operatorname{Id.inv}(X)$) au moyen des isomorphismes $I$ et $I'$ (resp. $\sh{J}_X$ et $I_X'$) précédents.
On notera que l'on a la relation
\[
\label{IV.21.2.7.1}
  D\geq0 \iff \sh{J}_X(D)\subset\sh{O}_X \quad\text{pour}\quad D\in\operatorname{Div}(X)
  \tag{21.2.7.1}
\]
comme il résulte aussitôt des définitions \sref{IV.21.1.6} et de \sref{IV.21.1.5.1};
en d'autres termes, l'image $\sh{J}_X(\operatorname{Div}^+(X))$ est l'ensemble des \emph{Idéaux} de $\sh{O}_X$ (dits aussi parfois \emph{Idéaux entiers}) qui sont des $\sh{O}_X$-Modules inversibles : un tel Idéal $\sh{J}$ est encore caractérisé par le fait que pour tout $x\in X$, il y a un voisinage ouvert $U$ de $x$ dans $X$ tel que $\sh{J}|U=\sh{O}_U.f$, où $f$ est un élément régulier de $\Gamma(U,\sh{O}_X)$.
L'ensemble $\sh{J}_X(\operatorname{Div}^+(X))$ de ces Idéaux est donc un sous-monoïde de $\operatorname{Id.inv}(X)$, égal à l'ensemble des éléments positifs pour une relation d'ordre compatible avec la structure de groupe de $\operatorname{Id.inv}(X)$, et il est immédiat que cette relation n'est autre que la relation \emph{opposée} à l'inclusion;
autrement dit, on a
\[
\label{IV.21.2.7.2}
  \sh{J}_1\subset\sh{J}_2 \iff \operatorname{div}(\sh{J}_1)\geq\operatorname{div}(\sh{J}_2)
  \tag{21.2.7.2}
\]
pour $\sh{J}_1$, $\sh{J}_2$ dans $\operatorname{Id.inv}(X)$.
\end{env}

\begin{env}[21.2.8]
\label{IV.21.2.8}
Pour tout diviseur $D$ sur $X$, on pose
\[
\label{IV.21.2.8.1}
  \sh{O}_X(D)=(\sh{J}_X(D))^{-1};
  \tag{21.2.8.1}
\]
$\sh{O}_X(D)$ est donc un Idéal fractionnaire inversible, défini de la façon suivante : pour tout ouvert $U$ de $X$ tel que $D|U=\operatorname{div}_U(f)$, où $f\in\Gamma(U,\sh{M}_X^*)$, $\sh{O}_X(D)|U$ est l'Idéal
